\documentclass[../../main.tex]{subfiles}
\begin{document}

\begin{center}
  {\fontsize{48pt}{58pt}\selectfont \textbf{2014}\normalsize\textbf{年度}}
\end{center}

\begin{table}[htb]
  \centering
  \caption{出題テーマと難易度}
  \begin{tabular}{|c|p{6.2cm}|c|c|c|}
    \hline
    \textbf{問題番号} & \textbf{テーマ} & \textbf{思考量} & \textbf{計算量} & \textbf{総合} \\
    \hline
    第1問 & 四角柱の平面断面の面積 & 易 & 易 & 易 \\
    \hline
    第2問 & 袋から球を取り出す確率と平均の極限 & 易 & 普 & 易 \\
    \hline
    第3問 & 二つの放物線の共有点と定積分 & 易 & 易 & 易 \\
    \hline
    第4問 & 指数関数で定める数列と不動点 & 普 & 普 & 普 \\
    \hline
    第5問 & 漸化式で定める数列の余り & 普 & 易 & 普 \\
    \hline
    第6問 & 動く線分が通過する平面領域 & 易 & 易 & 易 \\
    \hline
  \end{tabular}
\end{table}

\begin{table}[htb]
  \centering
  \caption{解答}
  \renewcommand{\arraystretch}{1.4}
  \begin{tabular}{|c|c|p{9cm}|}
    \hline
    \textbf{問題番号} & \textbf{小問} & \textbf{答え} \\
    \hline
    第1問 & (1) & $S=\sqrt{1+\tan^2\alpha+\tan^2\beta}$ \\
    \cline{2-3}
    & (2) & $\tan\alpha+\tan\beta=\dfrac{5}{6}$．
    $\alpha\le\beta$なら$\tan\alpha=\dfrac{1}{3}$ \\
    \hline
    第2問 & (1) & $p_1=\dfrac{1}{a+3}$，
    $p_2=\dfrac{a+2}{(a+3)(a+1)}$ \\
    \cline{2-3}
    & (2) & $p_n=\dfrac{1}{a+2}\left\{1-
    \dfrac{1}{a+3}\left(-\dfrac{1}{a+1}\right)^{n-1}\right\}$ \\
    \cline{2-3}
    & (3) & $\displaystyle\lim_{m\to\infty}\dfrac{1}{m}
    \sum_{n=1}^{m}p_n=\dfrac{1}{a+2}$ \\
    \hline
    第3問 & (1) & $a=-1-\sqrt{3}$，$b=-1+\sqrt{3}$ \\
    \cline{2-3}
    & (2) & $f(u)=\sqrt{-u^2-2u+2}\,(u^2+u+1)$ \\
    \cline{2-3}
    & (3) & $I=\dfrac{21\pi}{8}$ \\
    \hline
    第4問 & (1)〜(3) & 指定された関数の値域・数列の極限・不動点の存在を示す
    （証明問題） \\
    \hline
    第5問 & (1) & $b_{n+2}$は$b_{n+1}(b_n+1)$の$p$による余り
    （証明問題） \\
    \cline{2-3}
    & (2) & $b_1=b_4=b_7=b_{10}=2$，
    $b_2=b_5=b_8=3$，$b_3=b_6=b_9=9$ \\
    \cline{2-3}
    & (3) & 条件のもとで$b_n=b_m$（証明問題） \\
    \cline{2-3}
    & (4) & $a_2,a_3,\ldots$が$p$で割り切れなければ
    $a_1$も割り切れない（証明問題） \\
    \hline
    第6問 & (1) & $0\le s\le1$なら
    $\dfrac{\sqrt{3}}{3}(4-s)\le t\le
    \dfrac{\sqrt{3}}{6}(s^2+9)$．
    $1\le s\le2$なら$\sqrt{3}s\le t\le
    \dfrac{\sqrt{3}}{3}(s+4)$ \\
    \cline{2-3}
    & (2) & 領域$D$を図示（解答参照） \\
    \hline
  \end{tabular}
\end{table}

\paragraph{試験形態}
6問．試験時間は要確認．

\end{document}
